\chapter{Complexity — Machine}
%%% generated by the blueprint pipeline from TCSlib.Complexity.SpaceComplexity.SavitchARM.Machine
\section{Overview}

Declarations of the Lean module \texttt{TCSlib.Complexity.SpaceComplexity.SavitchARM.Machine}.

\section{Declarations}

\begin{definition}[Initial resource index]
\lean{Complexity.LogProg.SavitchARM.rn}
\label{Complexity.LogProg.SavitchARM.rn}
\leanok
Define $r_n$ to be the element $0$ of the finite index set $\{0,1,\ldots,17\}$.
\end{definition}

\begin{definition}[Acceptance resource index]
\lean{Complexity.LogProg.SavitchARM.ra}
\label{Complexity.LogProg.SavitchARM.ra}
\leanok
Define $r_a$ to be the element $1$ of the finite index set $\{0,1,\ldots,17\}$.
\end{definition}

\begin{definition}[Input resource index]
\lean{Complexity.LogProg.SavitchARM.rX}
\label{Complexity.LogProg.SavitchARM.rX}
\leanok
Define $r_X$ to be the element $2$ of the finite index set $\{0,1,\ldots,17\}$.
\end{definition}

\begin{definition}[Polynomial resource index]
\lean{Complexity.LogProg.SavitchARM.rP}
\label{Complexity.LogProg.SavitchARM.rP}
\leanok
Define $r_P$ to be the element $3$ of the finite index set $\{0,1,\ldots,17\}$.
\end{definition}

\begin{definition}[Time resource index]
\lean{Complexity.LogProg.SavitchARM.rT}
\label{Complexity.LogProg.SavitchARM.rT}
\leanok
Define $r_T$ to be the element $4$ of the finite index set $\{0,1,\ldots,17\}$.
\end{definition}

\begin{definition}[Index Five]
\lean{Complexity.LogProg.SavitchARM.rZ}
\label{Complexity.LogProg.SavitchARM.rZ}
\leanok
The distinguished index $r_Z$ is the element $5$ of the set of indices $\{0,1,\ldots,17\}$
(the 18-element type of indices).
\end{definition}

\begin{definition}[Index Six]
\lean{Complexity.LogProg.SavitchARM.rC}
\label{Complexity.LogProg.SavitchARM.rC}
\leanok
The distinguished index $r_C$ is the element $6$ of the set of indices $\{0,1,\ldots,17\}$
(the 18-element type of indices).
\end{definition}

\begin{definition}[Index Seven]
\lean{Complexity.LogProg.SavitchARM.rs}
\label{Complexity.LogProg.SavitchARM.rs}
\leanok
The distinguished index $r_s$ is the element $7$ of the set of indices $\{0,1,\ldots,17\}$
(the 18-element type of indices).
\end{definition}

\begin{definition}[Index Eight]
\lean{Complexity.LogProg.SavitchARM.rw}
\label{Complexity.LogProg.SavitchARM.rw}
\leanok
The distinguished index $r_w$ is the element $8$ of the set of indices $\{0,1,\ldots,17\}$
(the 18-element type of indices).
\end{definition}

\begin{definition}[Index Nine]
\lean{Complexity.LogProg.SavitchARM.rK}
\label{Complexity.LogProg.SavitchARM.rK}
\leanok
The distinguished index $r_K$ is the element $9$ of the set of indices $\{0,1,\ldots,17\}$
(the 18-element type of indices).
\end{definition}

\begin{definition}[Index Ten]
\lean{Complexity.LogProg.SavitchARM.rU}
\label{Complexity.LogProg.SavitchARM.rU}
\leanok
The distinguished index $r_U$ is the element $10$ of the set of indices $\{0,1,\ldots,17\}$
(the 18-element type of indices).
\end{definition}

\begin{definition}[Index Eleven]
\lean{Complexity.LogProg.SavitchARM.rV}
\label{Complexity.LogProg.SavitchARM.rV}
\leanok
The distinguished index $r_V$ is the element $11$ of the set of indices $\{0,1,\ldots,17\}$
(the 18-element type of indices).
\end{definition}

\begin{definition}[Index Twelve]
\lean{Complexity.LogProg.SavitchARM.rM}
\label{Complexity.LogProg.SavitchARM.rM}
\leanok
The distinguished index $r_M$ is the element $12$ of the set of indices $\{0,1,\ldots,17\}$
(the 18-element type of indices).
\end{definition}

\begin{definition}[Input register index]
\lean{Complexity.LogProg.SavitchARM.rI}
\label{Complexity.LogProg.SavitchARM.rI}
\leanok
The distinguished input-length register index is the element $13$ of the $18$-element
index set.
\end{definition}

\begin{definition}[Saved-endpoints register index]
\lean{Complexity.LogProg.SavitchARM.rSE}
\label{Complexity.LogProg.SavitchARM.rSE}
\leanok
The distinguished saved-endpoints register index is the element $14$ of the $18$-element
index set.
\end{definition}

\begin{definition}[Phase-stack register index]
\lean{Complexity.LogProg.SavitchARM.rSP}
\label{Complexity.LogProg.SavitchARM.rSP}
\leanok
The distinguished phase-stack register index is the element $15$ of the $18$-element
index set.
\end{definition}

\begin{definition}[Popped-endpoint register index]
\lean{Complexity.LogProg.SavitchARM.rE}
\label{Complexity.LogProg.SavitchARM.rE}
\leanok
The distinguished popped-endpoint register index is the element $16$ of the $18$-element
index set.
\end{definition}

\begin{definition}[Scratch register index]
\lean{Complexity.LogProg.SavitchARM.rQ}
\label{Complexity.LogProg.SavitchARM.rQ}
\leanok
The distinguished scratch-register index is the element $17$ of the $18$-element index
set.
\end{definition}

\begin{definition}[Machine register state]
\lean{Complexity.LogProg.SavitchARM.St}
\label{Complexity.LogProg.SavitchARM.St}
\leanok
A machine state is an assignment of a natural number to each of eighteen registers: the
input-length counter; the base $|x|+2$; the exponent counter and power accumulator; a
scratch register and the copy-and-shift scratch register; a loop counter; the window
radius; the code width (the recursion depth) and its bound $2^w$; the source and target
codes of the current node; its midpoint and level; the stack of saved endpoints, whose
digits have width $w$; the stack of phase bits; the last popped endpoint; and a scratch
register for the pop routine.
\end{definition}

\begin{definition}[Register vector of a state]
\lean{Complexity.LogProg.SavitchARM.St.toV}
\label{Complexity.LogProg.SavitchARM.St.toV}
\leanok
\uses{Complexity.LogProg.SavitchARM.St}
For each machine state, its register vector is the vector of length $18$ whose entries,
in order, are the input-length counter, the base, the exponent counter, the power
accumulator, the two scratch registers, the loop counter, the window radius, the code
width, the code bound, the source code, the target code, the midpoint, the level, the
saved-endpoints stack, the phase-bit stack, the last popped endpoint, and the
pop-routine scratch register.
\end{definition}

\begin{definition}[Zero machine state]
\lean{Complexity.LogProg.SavitchARM.St.zero}
\label{Complexity.LogProg.SavitchARM.St.zero}
\leanok
\uses{Complexity.LogProg.SavitchARM.St}
The zero machine state is the state in which all eighteen registers have value $0$.
\end{definition}

\begin{lemma}[Zero register vector]
\lean{Complexity.LogProg.SavitchARM.St.toV_zero}
\label{Complexity.LogProg.SavitchARM.St.toV_zero}
\leanok
\uses{Complexity.LogProg.SavitchARM.St.toV, Complexity.LogProg.SavitchARM.St.zero}
\difficulty{4}
The register vector of the all-zero state is the function that assigns zero to every
register index.
\end{lemma}

\begin{lemma}[Input counter coordinate]
\lean{Complexity.LogProg.SavitchARM.St.toV_n}
\label{Complexity.LogProg.SavitchARM.St.toV_n}
\leanok
\uses{Complexity.LogProg.SavitchARM.St, Complexity.LogProg.SavitchARM.St.toV, Complexity.LogProg.SavitchARM.rn}
\difficulty{1}
For any register state, the coordinate of its register vector at the input-counter index
is the state's input-length counter.
\end{lemma}

\begin{lemma}[Base register coordinate]
\lean{Complexity.LogProg.SavitchARM.St.toV_a}
\label{Complexity.LogProg.SavitchARM.St.toV_a}
\leanok
\uses{Complexity.LogProg.SavitchARM.St, Complexity.LogProg.SavitchARM.St.toV, Complexity.LogProg.SavitchARM.ra}
\difficulty{1}
For any register state, the coordinate of its register vector at the base-register index
is the state's base register.
\end{lemma}

\begin{lemma}[Exponent counter coordinate]
\lean{Complexity.LogProg.SavitchARM.St.toV_X}
\label{Complexity.LogProg.SavitchARM.St.toV_X}
\leanok
\uses{Complexity.LogProg.SavitchARM.St, Complexity.LogProg.SavitchARM.St.toV, Complexity.LogProg.SavitchARM.rX}
\difficulty{1}
For a state of the Savitch machine, the coordinate at index $2$ in its 18-entry register
vector equals the value of its exponent counter $X$.
\end{lemma}

\begin{lemma}[Power accumulator coordinate]
\lean{Complexity.LogProg.SavitchARM.St.toV_P}
\label{Complexity.LogProg.SavitchARM.St.toV_P}
\leanok
\uses{Complexity.LogProg.SavitchARM.St, Complexity.LogProg.SavitchARM.St.toV, Complexity.LogProg.SavitchARM.rP}
\difficulty{1}
For every machine state, the coordinate of its register vector indexed by the power
accumulator is equal to the value of its power accumulator register.
\end{lemma}

\begin{lemma}[Scratch register coordinate]
\lean{Complexity.LogProg.SavitchARM.St.toV_T}
\label{Complexity.LogProg.SavitchARM.St.toV_T}
\leanok
\uses{Complexity.LogProg.SavitchARM.St, Complexity.LogProg.SavitchARM.St.toV, Complexity.LogProg.SavitchARM.rT}
\difficulty{1}
For a state of the Savitch machine, the coordinate at index $4$ in its 18-entry register
vector equals the value of its scratch register $T$.
\end{lemma}

\begin{lemma}[Copy and shift scratch coordinate]
\lean{Complexity.LogProg.SavitchARM.St.toV_Z}
\label{Complexity.LogProg.SavitchARM.St.toV_Z}
\leanok
\uses{Complexity.LogProg.SavitchARM.St, Complexity.LogProg.SavitchARM.St.toV, Complexity.LogProg.SavitchARM.rZ}
\difficulty{1}
For a state of the Savitch machine, the coordinate at index $5$ in its 18-entry register
vector equals the value of its copy-and-shift-loop scratch register $Z$.
\end{lemma}

\begin{lemma}[Loop counter coordinate]
\lean{Complexity.LogProg.SavitchARM.St.toV_C}
\label{Complexity.LogProg.SavitchARM.St.toV_C}
\leanok
\uses{Complexity.LogProg.SavitchARM.St, Complexity.LogProg.SavitchARM.St.toV, Complexity.LogProg.SavitchARM.rC}
\difficulty{1}
For every machine state, the coordinate of its register vector indexed by the loop
counter is equal to the value of its loop counter register.
\end{lemma}

\begin{lemma}[Window radius coordinate]
\lean{Complexity.LogProg.SavitchARM.St.toV_s}
\label{Complexity.LogProg.SavitchARM.St.toV_s}
\leanok
\uses{Complexity.LogProg.SavitchARM.St, Complexity.LogProg.SavitchARM.St.toV, Complexity.LogProg.SavitchARM.rs}
\difficulty{1}
For any register state, the coordinate of its register vector at the window-radius index
is the state's window radius.
\end{lemma}

\begin{lemma}[Code width coordinate]
\lean{Complexity.LogProg.SavitchARM.St.toV_w}
\label{Complexity.LogProg.SavitchARM.St.toV_w}
\leanok
\uses{Complexity.LogProg.SavitchARM.St, Complexity.LogProg.SavitchARM.St.toV, Complexity.LogProg.SavitchARM.rw}
\difficulty{1}
For any register state, the coordinate of its register vector at the code-width index is
the state's code width.
\end{lemma}

\begin{lemma}[Code bound coordinate]
\lean{Complexity.LogProg.SavitchARM.St.toV_K}
\label{Complexity.LogProg.SavitchARM.St.toV_K}
\leanok
\uses{Complexity.LogProg.SavitchARM.St, Complexity.LogProg.SavitchARM.St.toV, Complexity.LogProg.SavitchARM.rK}
\difficulty{1}
For every machine state, the coordinate of its register vector indexed by the code bound
is equal to the value of its code bound register.
\end{lemma}

\begin{lemma}[Source code register coordinate]
\lean{Complexity.LogProg.SavitchARM.St.toV_U}
\label{Complexity.LogProg.SavitchARM.St.toV_U}
\leanok
\uses{Complexity.LogProg.SavitchARM.St, Complexity.LogProg.SavitchARM.St.toV, Complexity.LogProg.SavitchARM.rU}
\difficulty{1}
For a state of the Savitch machine, the coordinate at index $10$ in its 18-entry
register vector equals the value of its current-node source-code register $U$.
\end{lemma}

\begin{lemma}[Target code register coordinate]
\lean{Complexity.LogProg.SavitchARM.St.toV_V}
\label{Complexity.LogProg.SavitchARM.St.toV_V}
\leanok
\uses{Complexity.LogProg.SavitchARM.St, Complexity.LogProg.SavitchARM.St.toV, Complexity.LogProg.SavitchARM.rV}
\difficulty{1}
For a state of the Savitch machine, the coordinate at index $11$ in its 18-entry
register vector equals the value of its current-node target-code register $V$.
\end{lemma}

\begin{lemma}[Midpoint coordinate]
\lean{Complexity.LogProg.SavitchARM.St.toV_M}
\label{Complexity.LogProg.SavitchARM.St.toV_M}
\leanok
\uses{Complexity.LogProg.SavitchARM.St, Complexity.LogProg.SavitchARM.St.toV, Complexity.LogProg.SavitchARM.rM}
\difficulty{1}
For every machine state, the coordinate of its register vector indexed by the midpoint
is equal to the value of its midpoint register.
\end{lemma}

\begin{lemma}[Level coordinate]
\lean{Complexity.LogProg.SavitchARM.St.toV_I}
\label{Complexity.LogProg.SavitchARM.St.toV_I}
\leanok
\uses{Complexity.LogProg.SavitchARM.St, Complexity.LogProg.SavitchARM.St.toV, Complexity.LogProg.SavitchARM.rI}
\difficulty{1}
For every machine state, the coordinate of its register vector indexed by the level is
equal to the value of its level register.
\end{lemma}

\begin{lemma}[Saved endpoints register coordinate]
\lean{Complexity.LogProg.SavitchARM.St.toV_SE}
\label{Complexity.LogProg.SavitchARM.St.toV_SE}
\leanok
\uses{Complexity.LogProg.SavitchARM.St, Complexity.LogProg.SavitchARM.St.toV, Complexity.LogProg.SavitchARM.rSE}
\difficulty{1}
For a state of the Savitch machine, the coordinate at index $14$ in its 18-entry
register vector equals the value of its saved-endpoints stack register $SE$.
\end{lemma}

\begin{lemma}[Phase stack register coordinate]
\lean{Complexity.LogProg.SavitchARM.St.toV_SP}
\label{Complexity.LogProg.SavitchARM.St.toV_SP}
\leanok
\uses{Complexity.LogProg.SavitchARM.St, Complexity.LogProg.SavitchARM.St.toV, Complexity.LogProg.SavitchARM.rSP}
\difficulty{1}
For a state of the Savitch machine, the coordinate at index $15$ in its 18-entry
register vector equals the value of its phase-stack register $SP$.
\end{lemma}

\begin{lemma}[Popped endpoint coordinate]
\lean{Complexity.LogProg.SavitchARM.St.toV_E}
\label{Complexity.LogProg.SavitchARM.St.toV_E}
\leanok
\uses{Complexity.LogProg.SavitchARM.St, Complexity.LogProg.SavitchARM.St.toV, Complexity.LogProg.SavitchARM.rE}
\difficulty{1}
For every machine state, the coordinate of its register vector indexed by the popped
endpoint is equal to the value of its popped endpoint register.
\end{lemma}

\begin{lemma}[Target code register coordinate]
\lean{Complexity.LogProg.SavitchARM.St.toV_Q}
\label{Complexity.LogProg.SavitchARM.St.toV_Q}
\leanok
\uses{Complexity.LogProg.SavitchARM.St, Complexity.LogProg.SavitchARM.St.toV, Complexity.LogProg.SavitchARM.rQ}
\difficulty{1}
For a state of the Savitch machine, the coordinate at index $17$ in its 18-entry
register vector equals the value of its scratch register $Q$.
\end{lemma}

\begin{lemma}[Updating the input-length counter]
\lean{Complexity.LogProg.SavitchARM.St.upd_n}
\label{Complexity.LogProg.SavitchARM.St.upd_n}
\leanok
\uses{Complexity.LogProg.SavitchARM.St, Complexity.LogProg.SavitchARM.St.toV, Complexity.LogProg.SavitchARM.rn}
\difficulty{4}
For any register state $st$ and natural number $k$, replacing coordinate $0$ of the
vector of register values by $k$ gives the vector of register values of the state
obtained from $st$ by setting its input-length counter to $k$.
\end{lemma}

\begin{lemma}[Updating the base register]
\lean{Complexity.LogProg.SavitchARM.St.upd_a}
\label{Complexity.LogProg.SavitchARM.St.upd_a}
\leanok
\uses{Complexity.LogProg.SavitchARM.St, Complexity.LogProg.SavitchARM.St.toV, Complexity.LogProg.SavitchARM.ra}
\difficulty{4}
For any register state $st$ and natural number $k$, replacing coordinate $1$ of the
vector of register values by $k$ gives the vector of register values of the state
obtained from $st$ by setting its base register to $k$.
\end{lemma}

\begin{lemma}[Updating the exponent counter]
\lean{Complexity.LogProg.SavitchARM.St.upd_X}
\label{Complexity.LogProg.SavitchARM.St.upd_X}
\leanok
\uses{Complexity.LogProg.SavitchARM.St, Complexity.LogProg.SavitchARM.St.toV, Complexity.LogProg.SavitchARM.rX}
\difficulty{4}
For any register state $st$ and natural number $k$, replacing coordinate $2$ of the
vector of register values by $k$ gives the vector of register values of the state
obtained from $st$ by setting its exponent counter to $k$.
\end{lemma}

\begin{lemma}[Updating the power accumulator]
\lean{Complexity.LogProg.SavitchARM.St.upd_P}
\label{Complexity.LogProg.SavitchARM.St.upd_P}
\leanok
\uses{Complexity.LogProg.SavitchARM.St, Complexity.LogProg.SavitchARM.St.toV, Complexity.LogProg.SavitchARM.rP}
\difficulty{4}
For any state $st$ and natural number $k$, replacing coordinate $3$ of the register
vector of $st$ by $k$ gives the register vector of the state obtained by setting the
power accumulator $P$ to $k$.
\end{lemma}

\begin{lemma}[Updating the scratch register]
\lean{Complexity.LogProg.SavitchARM.St.upd_T}
\label{Complexity.LogProg.SavitchARM.St.upd_T}
\leanok
\uses{Complexity.LogProg.SavitchARM.St, Complexity.LogProg.SavitchARM.St.toV, Complexity.LogProg.SavitchARM.rT}
\difficulty{4}
For any state $st$ and natural number $k$, replacing coordinate $4$ of the register
vector of $st$ by $k$ gives the register vector of the state obtained by setting the
scratch register $T$ to $k$.
\end{lemma}

\begin{lemma}[Updating the copy-loop scratch register]
\lean{Complexity.LogProg.SavitchARM.St.upd_Z}
\label{Complexity.LogProg.SavitchARM.St.upd_Z}
\leanok
\uses{Complexity.LogProg.SavitchARM.St, Complexity.LogProg.SavitchARM.St.toV, Complexity.LogProg.SavitchARM.rZ}
\difficulty{4}
For any register state $st$ and natural number $k$, replacing coordinate $5$ of the
vector of register values by $k$ gives the vector of register values of the state
obtained from $st$ by setting its copy-and-shift scratch register to $k$.
\end{lemma}

\begin{lemma}[Updating the loop counter]
\lean{Complexity.LogProg.SavitchARM.St.upd_C}
\label{Complexity.LogProg.SavitchARM.St.upd_C}
\leanok
\uses{Complexity.LogProg.SavitchARM.St, Complexity.LogProg.SavitchARM.St.toV, Complexity.LogProg.SavitchARM.rC}
\difficulty{4}
For any register state and any natural number $k$, replacing the loop-counter coordinate
of its register vector by $k$ gives the register vector of the state with loop counter
set to $k$.
\end{lemma}

\begin{lemma}[Updating the window-radius register]
\lean{Complexity.LogProg.SavitchARM.St.upd_s}
\label{Complexity.LogProg.SavitchARM.St.upd_s}
\leanok
\uses{Complexity.LogProg.SavitchARM.St, Complexity.LogProg.SavitchARM.St.toV, Complexity.LogProg.SavitchARM.rs}
\difficulty{4}
For any register state $st$ and natural number $k$, replacing coordinate $7$ of the
vector of register values by $k$ gives the vector of register values of the state
obtained from $st$ by setting its window-radius register to $k$.
\end{lemma}

\begin{lemma}[Updating the code-width register]
\lean{Complexity.LogProg.SavitchARM.St.upd_w}
\label{Complexity.LogProg.SavitchARM.St.upd_w}
\leanok
\uses{Complexity.LogProg.SavitchARM.St, Complexity.LogProg.SavitchARM.St.toV, Complexity.LogProg.SavitchARM.rw}
\difficulty{4}
For any register state $st$ and natural number $k$, replacing coordinate $8$ of the
vector of register values by $k$ gives the vector of register values of the state
obtained from $st$ by setting its code-width register to $k$.
\end{lemma}

\begin{lemma}[Updating the code bound register]
\lean{Complexity.LogProg.SavitchARM.St.upd_K}
\label{Complexity.LogProg.SavitchARM.St.upd_K}
\leanok
\uses{Complexity.LogProg.SavitchARM.St, Complexity.LogProg.SavitchARM.St.toV, Complexity.LogProg.SavitchARM.rK}
\difficulty{4}
For any state $st$ and natural number $k$, replacing coordinate $9$ of the register
vector of $st$ by $k$ gives the register vector of the state obtained by setting the
code bound register $K$ to $k$.
\end{lemma}

\begin{lemma}[Updating the source code register]
\lean{Complexity.LogProg.SavitchARM.St.upd_U}
\label{Complexity.LogProg.SavitchARM.St.upd_U}
\leanok
\uses{Complexity.LogProg.SavitchARM.St, Complexity.LogProg.SavitchARM.St.toV, Complexity.LogProg.SavitchARM.rU}
\difficulty{4}
For any state $st$ and natural number $k$, replacing coordinate $10$ of the register
vector of $st$ by $k$ gives the register vector of the state obtained by setting the
source code register $U$ to $k$.
\end{lemma}

\begin{lemma}[Updating the target-code register]
\lean{Complexity.LogProg.SavitchARM.St.upd_V}
\label{Complexity.LogProg.SavitchARM.St.upd_V}
\leanok
\uses{Complexity.LogProg.SavitchARM.St, Complexity.LogProg.SavitchARM.St.toV, Complexity.LogProg.SavitchARM.rV}
\difficulty{4}
For any register state $st$ and natural number $k$, replacing coordinate $11$ of the
vector of register values by $k$ gives the vector of register values of the state
obtained from $st$ by setting its target-code register to $k$.
\end{lemma}

\begin{lemma}[Updating the midpoint register]
\lean{Complexity.LogProg.SavitchARM.St.upd_M}
\label{Complexity.LogProg.SavitchARM.St.upd_M}
\leanok
\uses{Complexity.LogProg.SavitchARM.St, Complexity.LogProg.SavitchARM.St.toV, Complexity.LogProg.SavitchARM.rM}
\difficulty{4}
For any state $st$ and natural number $k$, replacing coordinate $12$ of the register
vector of $st$ by $k$ gives the register vector of the state obtained by setting the
midpoint register $M$ to $k$.
\end{lemma}

\begin{lemma}[Updating the level counter]
\lean{Complexity.LogProg.SavitchARM.St.upd_I}
\label{Complexity.LogProg.SavitchARM.St.upd_I}
\leanok
\uses{Complexity.LogProg.SavitchARM.St, Complexity.LogProg.SavitchARM.St.toV, Complexity.LogProg.SavitchARM.rI}
\difficulty{4}
For any register state and any natural number $k$, replacing the current-level
coordinate of its register vector by $k$ gives the register vector of the state with
current level set to $k$.
\end{lemma}

\begin{lemma}[Updating the saved endpoints stack]
\lean{Complexity.LogProg.SavitchARM.St.upd_SE}
\label{Complexity.LogProg.SavitchARM.St.upd_SE}
\leanok
\uses{Complexity.LogProg.SavitchARM.St, Complexity.LogProg.SavitchARM.St.toV, Complexity.LogProg.SavitchARM.rSE}
\difficulty{4}
For any state $st$ and natural number $k$, replacing coordinate $14$ of the register
vector of $st$ by $k$ gives the register vector of the state obtained by setting the
saved endpoints stack $SE$ to $k$.
\end{lemma}

\begin{lemma}[Updating the phase bits stack]
\lean{Complexity.LogProg.SavitchARM.St.upd_SP}
\label{Complexity.LogProg.SavitchARM.St.upd_SP}
\leanok
\uses{Complexity.LogProg.SavitchARM.St, Complexity.LogProg.SavitchARM.St.toV, Complexity.LogProg.SavitchARM.rSP}
\difficulty{4}
For any state $st$ and natural number $k$, replacing coordinate $15$ of the register
vector of $st$ by $k$ gives the register vector of the state obtained by setting the
phase bits stack $SP$ to $k$.
\end{lemma}

\begin{lemma}[Updating the last endpoint]
\lean{Complexity.LogProg.SavitchARM.St.upd_E}
\label{Complexity.LogProg.SavitchARM.St.upd_E}
\leanok
\uses{Complexity.LogProg.SavitchARM.St, Complexity.LogProg.SavitchARM.St.toV, Complexity.LogProg.SavitchARM.rE}
\difficulty{4}
For any register state and any natural number $k$, replacing the last-endpoint
coordinate of its register vector by $k$ gives the register vector of the state with
last endpoint set to $k$.
\end{lemma}

\begin{lemma}[Updating the pop routine scratch register]
\lean{Complexity.LogProg.SavitchARM.St.upd_Q}
\label{Complexity.LogProg.SavitchARM.St.upd_Q}
\leanok
\uses{Complexity.LogProg.SavitchARM.St, Complexity.LogProg.SavitchARM.St.toV, Complexity.LogProg.SavitchARM.rQ}
\difficulty{3}
For any state $st$ and natural number $k$, replacing coordinate $17$ of the register
vector of $st$ by $k$ gives the register vector of the state obtained by setting the pop
routine scratch register $Q$ to $k$.
\end{lemma}

\begin{definition}[Uniform register bound]
\lean{Complexity.LogProg.SavitchARM.St.Bd}
\label{Complexity.LogProg.SavitchARM.St.Bd}
\leanok
\uses{Complexity.LogProg.SavitchARM.St}
For a natural number $B$ and a machine state with counters and registers
$n,a,X,P,T,Z,C,s,w,K,U,V,M,I,SE,SP,E,Q$, the state is bounded by $B$ exactly when each
of these eighteen quantities is at most $B$.
\end{definition}

\begin{lemma}[Register vector respects the bound]
\lean{Complexity.LogProg.SavitchARM.St.Bd.toV}
\label{Complexity.LogProg.SavitchARM.St.Bd.toV}
\leanok
\uses{Complexity.LogProg.SavitchARM.St, Complexity.LogProg.SavitchARM.St.Bd, Complexity.LogProg.SavitchARM.St.toV}
\difficulty{4}
Let $B$ be a natural number and let a machine state have all eighteen registers at most
$B$. Then every coordinate of its register vector is at most $B$.
\end{lemma}

\begin{definition}[Move-loop sites]
\lean{Complexity.LogProg.SavitchARM.MvS}
\label{Complexity.LogProg.SavitchARM.MvS}
\leanok
Let a move-loop site be one of the following finitely many operations: copy $n$ into
$a$; copy a product from $T$ into $P$, with a Boolean parameter; copy $P$ into $w$; copy
either $U$ or $V$ into the endpoint stack, according to a Boolean parameter; copy the
doubled phase stack from $T$ into $SP$, with a Boolean parameter; copy $M$ into $V$;
copy $V$ into $U$; copy the popped endpoint $E$ into $V$; move the last endpoint-stack
digit into $E$, with a Boolean parameter; copy the remaining stack from $T$ into $SE$,
with a Boolean parameter; copy $V$ into $M$; copy $E$ into $V$; copy $E$ into $U$; copy
$U$ into $M$; or copy $E$ into $U$.
\end{definition}

\begin{definition}[Copy-loop sites]
\lean{Complexity.LogProg.SavitchARM.CpS}
\label{Complexity.LogProg.SavitchARM.CpS}
\leanok
Let a copy-loop site be one of the following finitely many operations: choose the
multiplier, either $P$ or $a$, according to a Boolean parameter, and copy it into
counter $C$; copy $P$ into accumulator $T$, with a Boolean parameter; copy $P$ into
radius $s$; copy $w$ into $C$ for the shift count used to form $K=2^w$; copy $w$ into
level $I$; copy $w$ into $C$ for shifting the endpoint stack, with a Boolean parameter;
copy $SE$ into $T$ for popping, with a Boolean parameter; copy $w$ into $C$ for the
stack-halving count, with a Boolean parameter; copy the halved stack from $T$ into $Q$,
with a Boolean parameter; or copy $w$ into $C$ for shifting $Q$, with a Boolean
parameter.
\end{definition}

\begin{definition}[Shift-loop sites]
\lean{Complexity.LogProg.SavitchARM.ShS}
\label{Complexity.LogProg.SavitchARM.ShS}
\leanok
A shift-loop site is one of three operations: multiply $K$ by $2^w$; multiply the
endpoint stack $SE$ by $2^w$ to make room for a digit, with a Boolean parameter; or
multiply $Q$ by $2^w$ to shift back the stack with its last digit removed, with a
Boolean parameter.
\end{definition}

\begin{definition}[Machine-program labels]
\lean{Complexity.LogProg.SavitchARM.Lb}
\label{Complexity.LogProg.SavitchARM.Lb}
\leanok
\uses{Complexity.LogProg.SavitchARM.CpS, Complexity.LogProg.SavitchARM.MvS, Complexity.LogProg.SavitchARM.ShS}
For a natural number $E$, let the program-label type consist of the labels for the
setup, search, recursion, return, and stack operations: input scanning and counter
updates; power and multiplication loops; initialization and outer-search tests and
queries; root, accept, and reject operations; recursive calls and base-case tests;
midpoint iteration and phase handling; and the move, doubling, subtraction, halving,
copy, and shift loops. The label for the $j$th exponent-counter increment is
parameterized by $j$ in the range $0\le j\le E$, and the move, doubling, subtraction,
halving, copy, and shift loop labels carry, respectively, a move site and phase in
$\{0,1,2\}$, a Boolean and phase in $\{0,1,2,3\}$, a Boolean and phase in $\{0,1,2\}$, a
Boolean and phase in $\{0,1,2\}$, a copy site and phase in $\{0,1,\ldots,6\}$, or a
shift site and phase in $\{0,1,\ldots,8\}$.
\end{definition}

\begin{definition}[Source register of each move]
\lean{Complexity.LogProg.SavitchARM.mvSrc}
\label{Complexity.LogProg.SavitchARM.mvSrc}
\leanok
\uses{Complexity.LogProg.SavitchARM.MvS, Complexity.LogProg.SavitchARM.rE, Complexity.LogProg.SavitchARM.rM, Complexity.LogProg.SavitchARM.rP, Complexity.LogProg.SavitchARM.rSE, Complexity.LogProg.SavitchARM.rT, Complexity.LogProg.SavitchARM.rU, Complexity.LogProg.SavitchARM.rV, Complexity.LogProg.SavitchARM.rn}
For a move site, define its source-register index in a register file indexed from $0$ to
$17$ as follows: the source indices for copying $n$ to $a$, $T$ to $P$, $P$ to $w$, the
selected endpoint to $SE$, $T$ to $SP$, $M$ to $V$, $V$ to $U$, $E$ to $V$, $SE$ to $E$,
$T$ to $SE$, $V$ to $M$, $E$ to $V$, $E$ to $U$, $U$ to $M$, and $E$ to $U$ are,
respectively, $0,4,3,10$ or $11$ (according to the Boolean choice),
$4,12,11,16,14,4,11,16,16,10,16$.
\end{definition}

\begin{definition}[Destination register of each move]
\lean{Complexity.LogProg.SavitchARM.mvDst}
\label{Complexity.LogProg.SavitchARM.mvDst}
\leanok
\uses{Complexity.LogProg.SavitchARM.MvS, Complexity.LogProg.SavitchARM.rE, Complexity.LogProg.SavitchARM.rM, Complexity.LogProg.SavitchARM.rP, Complexity.LogProg.SavitchARM.rSE, Complexity.LogProg.SavitchARM.rSP, Complexity.LogProg.SavitchARM.rU, Complexity.LogProg.SavitchARM.rV, Complexity.LogProg.SavitchARM.ra, Complexity.LogProg.SavitchARM.rw}
For a move site, define its destination-register index in a register file indexed from
$0$ to $17$ as follows: the destination indices for copying $n$ to $a$, $T$ to $P$, $P$
to $w$, the selected endpoint to $SE$, $T$ to $SP$, $M$ to $V$, $V$ to $U$, $E$ to $V$,
$SE$ to $E$, $T$ to $SE$, $V$ to $M$, $E$ to $V$, $E$ to $U$, $U$ to $M$, and $E$ to $U$
are, respectively, $1,3,8,14,15,11,10,11,16,14,12,11,10,12,10$.
\end{definition}

\begin{definition}[Exit label of each move loop]
\lean{Complexity.LogProg.SavitchARM.mvExit}
\label{Complexity.LogProg.SavitchARM.mvExit}
\leanok
\uses{Complexity.LogProg.SavitchARM.Lb, Complexity.LogProg.SavitchARM.MvS}
For each move site, define its exit program label as follows: copying $n$ into $a$ exits
to the first increment of $a$; copying $T$ into $P$ exits to the $w$-copy loop when the
Boolean parameter is true and otherwise to the power test; copying $P$ into $w$ exits to
clearing $a$; pushing an endpoint into $SE$ exits to the corresponding doubling loop;
pushing $T$ into $SP$ exits to incrementing the phase when its Boolean parameter is true
and otherwise to the descent loop; copying $M$ into $V$ exits to decrementing the level;
the two phase-$0$ success copies exit successively to the second success copy and then
to the level decrement; popping the last endpoint digit into $E$ exits to copying the
remaining stack into $SE$; that remaining-stack copy exits to restoring the level; the
phase-$0$ failure copies exit to the next midpoint and then to the midpoint loop; the
phase-$1$ success copy exits to the root-return test; and the phase-$1$ failure copies
exit to the next midpoint and then to the midpoint loop. The exit labels for the stack
push and pop operations retain their Boolean parameter, and the level-restoration label
also retains its Boolean parameter.
\end{definition}

\begin{definition}[Source registers for copy operations]
\lean{Complexity.LogProg.SavitchARM.cpSrc}
\label{Complexity.LogProg.SavitchARM.cpSrc}
\leanok
\uses{Complexity.LogProg.SavitchARM.CpS, Complexity.LogProg.SavitchARM.rP, Complexity.LogProg.SavitchARM.rSE, Complexity.LogProg.SavitchARM.rT, Complexity.LogProg.SavitchARM.ra, Complexity.LogProg.SavitchARM.rw}
For each copy-operation site, assign a source register among the 18 registers as
follows: multiplication into the counter uses the power register when the Boolean
parameter is true and the base register otherwise; multiplication into the accumulator
and setup of the radius use the power register; setup, root, push, and both pop-count
operations use the width register; popping the stack into the accumulator uses the stack
register, and moving the halved stack into the quotient register uses the accumulator
register.
\end{definition}

\begin{definition}[Destination registers for copy operations]
\lean{Complexity.LogProg.SavitchARM.cpDst}
\label{Complexity.LogProg.SavitchARM.cpDst}
\leanok
\uses{Complexity.LogProg.SavitchARM.CpS, Complexity.LogProg.SavitchARM.rC, Complexity.LogProg.SavitchARM.rI, Complexity.LogProg.SavitchARM.rQ, Complexity.LogProg.SavitchARM.rT, Complexity.LogProg.SavitchARM.rs}
For each copy-operation site, assign its destination among the 18 registers as follows:
multiplication into the counter, setup of the shift value, push, and both pop-count
operations target the counter register; multiplication into the accumulator and popping
the stack into the accumulator target the accumulator register; setup of the radius
targets the radius register; the root operation targets the level register; and moving
the halved stack into the quotient register targets the quotient register.
\end{definition}

\begin{definition}[Exit labels for copy operations]
\lean{Complexity.LogProg.SavitchARM.cpExit}
\label{Complexity.LogProg.SavitchARM.cpExit}
\leanok
\uses{Complexity.LogProg.SavitchARM.CpS, Complexity.LogProg.SavitchARM.Lb}
For each copy-operation site, assign its continuation label as follows: either
multiplication site continues at the corresponding multiplication test; radius setup
continues at phase $0$ of the squaring multiplication site; shift-value setup continues
at phase $0$ of the shift of the setup register; root-level copying continues at the
recursive-call label; pushing continues at phase $0$ of the corresponding push shift;
copying the stack during a pop continues at phase $0$ of the first pop-count copy; the
first pop-count copy continues at phase $0$ of the corresponding halving operation;
copying the halved stack continues at phase $0$ of the second pop-count copy; and the
second pop-count copy continues at phase $0$ of the corresponding pop shift.
\end{definition}

\begin{definition}[Registers used by shift operations]
\lean{Complexity.LogProg.SavitchARM.shReg}
\label{Complexity.LogProg.SavitchARM.shReg}
\leanok
\uses{Complexity.LogProg.SavitchARM.ShS, Complexity.LogProg.SavitchARM.rK, Complexity.LogProg.SavitchARM.rQ, Complexity.LogProg.SavitchARM.rSE}
For each shift-operation site, assign its register among the 18 registers: shifting the
setup value uses the setup register, pushing uses the stack register, and popping uses
the quotient register.
\end{definition}

\begin{definition}[Exit labels for shift operations]
\lean{Complexity.LogProg.SavitchARM.shExit}
\label{Complexity.LogProg.SavitchARM.shExit}
\leanok
\uses{Complexity.LogProg.SavitchARM.Lb, Complexity.LogProg.SavitchARM.ShS}
For each shift-operation site, assign its continuation label as follows: shifting the
setup value continues at the outer search label; pushing continues at phase $0$ of the
corresponding saved-endpoint move; and popping continues at phase $0$ of the
corresponding subtraction operation. The label type is parameterized by a natural number
$E$.
\end{definition}

\begin{definition}[Instruction phases for a register move]
\lean{Complexity.LogProg.SavitchARM.mvI}
\label{Complexity.LogProg.SavitchARM.mvI}
\leanok
\uses{Complexity.LogProg.Ins, Complexity.LogProg.SavitchARM.Lb}
Let $E$ be a natural number, let $s$ and $d$ be registers among 18 registers, let $x$ be
an exit label, let $m$ assign a label to each of three phases, and let $j$ be one of
those phases. The instruction at phase $j$ is a zero test on register $s$, branching to
$x$ if it is zero and to phase $1$ otherwise when $j=0$; it decrements $s$ and continues
to phase $2$ when $j=1$; and in the remaining phase it increments $d$ and returns to
phase $0$.
\end{definition}

\begin{definition}[Instruction phases for a doubling move]
\lean{Complexity.LogProg.SavitchARM.dmI}
\label{Complexity.LogProg.SavitchARM.dmI}
\leanok
\uses{Complexity.LogProg.Ins, Complexity.LogProg.SavitchARM.Lb}
Let $E$ be a natural number, let $s$ and $d$ be registers among 18 registers, let $x$ be
an exit label, let $m$ assign a label to each of four phases, and let $j$ be one of
those phases. At phase $0$, test whether $s$ is zero, branching to $x$ if so and to
phase $1$ otherwise; at phase $1$, decrement $s$ and continue to phase $2$; at phase
$2$, increment $d$ and continue to phase $3$; and at phase $3$, increment $d$ and return
to phase $0$.
\end{definition}

\begin{definition}[Instruction phases for subtraction]
\lean{Complexity.LogProg.SavitchARM.sbI}
\label{Complexity.LogProg.SavitchARM.sbI}
\leanok
\uses{Complexity.LogProg.Ins, Complexity.LogProg.SavitchARM.Lb}
Let $E$ be a natural number, let $q$ and $r$ be registers among 18 registers, let $x$ be
an exit label, let $m$ assign a label to each of three phases, and let $j$ be one of
those phases. At phase $0$, test whether $q$ is zero, branching to $x$ if so and to
phase $1$ otherwise; at phase $1$, decrement $q$ and continue to phase $2$; and in the
remaining phase, decrement $r$ and return to phase $0$.
\end{definition}

\begin{definition}[Halving-loop instruction]
\lean{Complexity.LogProg.SavitchARM.hvI}
\label{Complexity.LogProg.SavitchARM.hvI}
\leanok
\uses{Complexity.LogProg.Ins, Complexity.LogProg.SavitchARM.Lb}
For a natural number $E$, two registers $c,r
ame$ in a set of $18$ registers, an exit label $e$, a function $m$ assigning one of the
labels $m(0),m(1),m(2)$ to each phase, and a phase $j\in\{0,1,2\}$, define the
instruction as follows: at phase $0$, branch on whether register $c$ is zero, going to
$e$ if it is and to $m(1)$ otherwise; at phase $1$, decrement $c$ and go to $m(2)$; at
phase $2$, halve register $r$ and go to $m(0)$.
\end{definition}

\begin{definition}[Copy-loop instruction]
\lean{Complexity.LogProg.SavitchARM.cpI}
\label{Complexity.LogProg.SavitchARM.cpI}
\leanok
\uses{Complexity.LogProg.Ins, Complexity.LogProg.SavitchARM.Lb, Complexity.LogProg.SavitchARM.rZ}
For a natural number $E$, source and destination registers $s,d$ among $18$ registers,
an exit label $e$, a function $m$ assigning a label to each of seven phases, and a phase
$j\in\{0,1,\ldots,6\}$, define the instruction by these cases: at phase $0$, branch on
whether $s=0$, going to $m(4)$ if so and to $m(1)$ otherwise; at phase $1$, decrement
$s$ and go to $m(2)$; at phase $2$, increment $d$ and go to $m(3)$; at phase $3$,
increment a scratch register $z$ and go to $m(0)$; at phase $4$, branch on whether
$z=0$, going to $e$ if so and to $m(5)$ otherwise; at phase $5$, decrement $z$ and go to
$m(6)$; at phase $6$, increment $s$ and go to $m(4)$.
\end{definition}

\begin{definition}[Shift-loop instruction]
\lean{Complexity.LogProg.SavitchARM.shI}
\label{Complexity.LogProg.SavitchARM.shI}
\leanok
\uses{Complexity.LogProg.Ins, Complexity.LogProg.SavitchARM.Lb, Complexity.LogProg.SavitchARM.rC, Complexity.LogProg.SavitchARM.rZ}
For a natural number $E$, a register $r$ among $18$ registers, an exit label $e$, a
function $m$ assigning a label to each of nine phases, and a phase
$j\in\{0,1,\ldots,8\}$, define the instruction by these cases: at phase $0$, branch on
whether the loop counter $c$ is zero, going to $e$ if so and to $m(1)$ otherwise; at
phase $1$, decrement $c$ and go to $m(2)$; at phase $2$, branch on whether $r=0$, going
to $m(6)$ if so and to $m(3)$ otherwise; at phase $3$, decrement $r$ and go to $m(4)$;
at phase $4$, increment the scratch register $z$ and go to $m(5)$; at phase $5$,
increment $z$ and go to $m(2)$; at phase $6$, branch on whether $z=0$, going to $m(0)$
if so and to $m(7)$ otherwise; at phase $7$, decrement $z$ and go to $m(8)$; at phase
$8$, increment $r$ and go to $m(6)$.
\end{definition}

\begin{definition}[Savitch machine program]
\lean{Complexity.LogProg.SavitchARM.A}
\label{Complexity.LogProg.SavitchARM.A}
\leanok
\uses{Complexity.LogProg.ARM, Complexity.LogProg.SavitchARM.Lb, Complexity.LogProg.SavitchARM.cpDst, Complexity.LogProg.SavitchARM.cpExit, Complexity.LogProg.SavitchARM.cpI, Complexity.LogProg.SavitchARM.cpSrc, Complexity.LogProg.SavitchARM.dmI, Complexity.LogProg.SavitchARM.hvI, Complexity.LogProg.SavitchARM.mvDst, Complexity.LogProg.SavitchARM.mvExit, Complexity.LogProg.SavitchARM.mvI, Complexity.LogProg.SavitchARM.mvSrc, Complexity.LogProg.SavitchARM.rC, Complexity.LogProg.SavitchARM.rI, Complexity.LogProg.SavitchARM.rK, Complexity.LogProg.SavitchARM.rM, Complexity.LogProg.SavitchARM.rP, Complexity.LogProg.SavitchARM.rQ, Complexity.LogProg.SavitchARM.rSE, Complexity.LogProg.SavitchARM.rSP, Complexity.LogProg.SavitchARM.rT, Complexity.LogProg.SavitchARM.rU, Complexity.LogProg.SavitchARM.rV, Complexity.LogProg.SavitchARM.rX, Complexity.LogProg.SavitchARM.ra, Complexity.LogProg.SavitchARM.rn, Complexity.LogProg.SavitchARM.rs, Complexity.LogProg.SavitchARM.rw, Complexity.LogProg.SavitchARM.sbI, Complexity.LogProg.SavitchARM.shExit, Complexity.LogProg.SavitchARM.shI, Complexity.LogProg.SavitchARM.shReg}
For a natural number $E$, define a register-machine program with $18$ registers, four
deciders, and labels in the family indexed by $E$. At each label it supplies the
corresponding instruction: the program scans the input length, computes its parameters
and code bound, searches pairs of codes, and performs the recursive reachability
procedure using the move, copy, shift, doubling, subtraction, and halving loops
specified by the labels.
\end{definition}
